\subsection{APPLICATION TO ASYMPTOTIC EXPANSIONS}

The Euler-Maclaurin formula allows one to give a more complete solution (in the most important cases) to the problem treated in V, p.~238 to 242, that of obtaining an asymptotic expansion of the partial sum \(s_n = \sum_{m=0}^{n} g(m)\) (resp. of the remainder \(r_n = \sum_{m=n+1}^{\infty} g(m)\) ), where \(g\) is a scalar function, \(> 0\) , monotone, defined on \([0, +\infty[\) .

We shall restrict ourselves to the case where g is an (H) function (V, p.~252), of order 0 relative to \(e^{x}\) ; in other words, one has the relation \(g' \ll g\) ; from this relation one deduces \(g^{(k+1)} \ll g^{(k)}\) for every integer k \textgreater{} 0 such that none of the derivatives \(g^{(h)}\) of order \(h \leqslant k\) is equivalent to a constant (V, p.~232, prop. 7). Let p be an integer such that none of the derivatives \(g^{(h)}\) of order \(h \leqslant 2p\) is equivalent to a constant. First we assume that the series with general term \(g(n)\) has infinite sum, and distinguish several cases:

\(1^{\circ}\left|g^{(2p - 1)}(n)\right|\) tends to \(+\infty\) with \(n\) ; we have the same, by the hypothesis, for \(\left|g^{(2k - 1)}(n)\right|\) for \(1\leqslant k\leqslant p\) ; further, since \(g^{(2p + 1)}\) is monotone on a neighbourhood of \(+\infty\) , the formula (4) of VI, p.~289, gives \(\mathrm{T}_p(0,n) = O\big(g^{(2p)}(n + 1)\big) = o\big(g^{(2p - 1)}(n + 1)\big)\) ; the Euler-Maclaurin formula, applied for \(x = 0\) , shows that

\[
\begin{array}{l} s _ {n} = \sum_ {m = 0} ^ {n} g (m) = \int_ {0} ^ {n + 1} g (t) d t - \frac {1}{2} g (n + 1) \\ \qquad + \sum_ {k = 1} ^ {p} \frac {b _ {2 k}}{(2 k) !} g ^ {(2 k - 1)} (n + 1) + o \big (g ^ {(2 p - 1)} (n + 1) \big) \end{array}
\]

each of the terms of this sum being negligible relative to the preceding one; on expanding each of them relative to a comparison scale E one will then have an asymptotic expansion for \(s_{n}\) .

\(2^{\circ}\) Now suppose that for an index \(q\) such that \(1 \leqslant q \leqslant p\) we have \(\left|g^{(2q-1)}(n)\right|\) tending to \(+\infty\) with \(n\) , but that \(g^{(2k-1)}(n)\) tends to 0 for \(k > q\) . Since \(g^{(2p+1)}\) is monotone on a neighbourhood of \(+\infty\) the integral \(\int_0^\infty \left|g^{(2p+1)}(u)\right| du\) converges, and one can then write

\[
\begin{array}{l} s _ {n} = \sum_ {m = 0} ^ {n} g (m) = \int_ {0} ^ {n + 1} g (t) d t - \frac {1}{2} g (n + 1) + \sum_ {k = 1} ^ {q} \frac {b _ {2 k}}{(2 k) !} g ^ {(2 k - 1)} (n + 1) + C \\ \qquad + \sum_ {k = q + 1} ^ {p} \frac {b _ {2 k}}{(2 k) !} g ^ {(2 k - 1)} (n + 1) + o \big (g ^ {(2 p - 1)} (n + 1) \big) \end{array}
\]

where C is a constant: indeed

\[
\int_ {n + 1} ^ {\infty} \left| g ^ {(2 p + 1)} (u) \right| d u = O \big (g ^ {(2 p)} (n + 1) \big) = o \big (g ^ {(2 p - 1)} (n + 1) \big).
\]

The same formula is valid when \(g(n)\) itself tends to 0. Finally, when the series with general term \(g(n)\) converges one has, for the remainder \(r_n = \sum_{m=n+1}^{\infty} g(m)\) , the expansion

\[
\begin{array}{l} r _ {n} = \sum_ {m = n + 1} ^ {\infty} g (m) = \int_ {n + 1} ^ {\infty} g (t) d t + \frac {1}{2} g (n + 1) \\ \qquad - \sum_ {k = 1} ^ {p} \frac {b _ {2 k}}{(2 k) !} g ^ {(2 k - 1)} (n + 1) + o \big (g ^ {(2 p - 1)} (n + 1) \big). \end{array}
\]

\setcounter{exercisegroup}{0}
\refstepcounter{exercisegroup}
